% Optional component-study subsection; does not modify the frozen V1 manuscript.
\subsection{Fragment-repair component and shared-head integration screens}
We evaluated established selective fragment recovery with block acknowledgements
as a stronger MAC comparator, rather than as a new ARQ protocol. Each 500-byte
datagram traversed two hops in six fragments. Paid reports, grants, cumulative
bitmaps and fixed radio windows were modeled explicitly. A source retained custody
until receiving a complete delivery receipt. These are synthetic simulations,
not hardware measurements or an RFC-conformant implementation.

In a corrected single-source development screen (64 paired seeds, 96 scenario
cells), block acknowledgement increased timely delivery from 11.188\% to
14.458\% and energy efficiency from 33.8514 to 43.2781 packets/J relative to
whole-packet retry. The corresponding relative improvements were 29.226\% and
27.847\%. Two additional proposed extensions failed their prespecified gates:
affordable block prefixes and a last-frame completion guard. An earlier experiment
was invalidated because variable transaction timing undercharged ambiguous waits;
its performance is excluded.

A separate shared-head integration screen used 32 fresh paired development seeds,
96 scenario cells and three policies (9216 episodes). One, four or twelve sources
shared a head battery, two live reassembly buffers and a one-second serialized
frame. Timely delivery was 12.558\%, 13.493\% and 17.121\% for whole-packet retry,
per-fragment selective retry and block acknowledgement, respectively; efficiency
was 41.4527, 44.5781 and 57.8723 packets/J. Relative to whole-packet retry, the
block scheme improved delivery by 36.340\% (paired 95\% interval
[33.499\%,39.197\%]) and efficiency by 39.611\% ([36.695\%,42.565\%]).
Against selective retry, gains were 26.892\% and 29.822\%. All four
family-adjusted bootstrap lower bounds exceeded the prespecified 1\% margin.
Resampling was over seeds, retaining scenario and policy pairing; percentile
coverage is approximate. Results are development screens, not confirmation.

The low absolute delivery, reliable-control assumption, conservative radio
envelope and short horizon limit applicability. Heads were scripted and dedicated;
this study did not reproduce the full mobile HEART-CH workload. Consequently,
these percentages do not quantify improvement over HTA-MAC V1 or other published
systems, and do not certify the original attribution or network-performance
gates. No neural policy was trained. Full-network integration and independent
confirmation remain necessary.
